\section{Scope, Method, and Qualifications}

\subsection*{Subject, Series Context, and Governing Question}

\emph{The Due Return} is Part~II of \emph{The Creature Before God}, a discursive Catholic theological series examining creaturely dependence, worship, and divine friendship. It addresses a single governing question: if existence was never owed to the creature and can never be repaid in kind, in what sense is worship owed in justice, how does that due return take embodied, liturgical, moral, and merciful form even before affection matures, and why can no fulfilled act close the creature's account?

The intended reader is a believer or serious inquirer seeking a primary-source account of Christian worship, moral duty, sacrifice, reconciliation, and freedom rooted in Sacred Scripture, the Fathers, St.\ Thomas Aquinas, ecumenical councils, and the Church's pastoral and spiritual tradition.

\subsection*{Reason, Revelation, and Authority Classification}

The argument maintains a clear boundary between what natural reason can discern and what belongs solely to divine revelation and Christian faith:

\begin{itemize}
\item \textbf{Natural reason and philosophical theology} establish radical creaturely dependence, the impossibility of strict equivalence with the Creator, the generic moral duty of gratitude and outward signs of subjection to God alone, and the necessity of determinate positive norms for common life (developed here chiefly through Aquinas, \emph{ST} I-II, qq.\ 90--91, 94; II-II, qq.\ 58, 80--81, 85, 106).
\item \textbf{Divine revelation and Catholic dogma} establish the historical economy of Exodus and covenant, the Incarnation and two natural wills of Christ (Constantinople~III), Christ's unique saving sacrifice represented in the Mass (Council of Trent, Session~XXII), gratuitous justification and grace-grounded merit (Trent, Session~VI), the sacrament of Penance and Christ-centered satisfaction (Trent, Session~XIV), and the obligation and freedom of supernatural faith (Vatican~I, \emph{Dei Filius}, ch.\ 3 and Canon~III.5).
\item \textbf{Authoritative non-definitive magisterium} supplies the Church's pastoral and doctrinal exposition of active liturgical participation (\emph{Sacrosanctum Concilium}), the common priesthood of the faithful (\emph{Lumen gentium}), immunity from external coercion in religious acts (\emph{Dignitatis humanae}), Sunday as the Day of the Church (\emph{Dies Domini}), and catechetical synthesis (\emph{Catechism of the Catholic Church} and the 1566 \emph{Roman Catechism}).
\item \textbf{Patristic, scholastic, and saintly witnesses} (St.\ Justin Martyr, St.\ Irenaeus, St.\ Gregory of Nazianzus, St.\ John Chrysostom, St.\ Augustine, St.\ Benedict, St.\ Thomas Aquinas, and St.\ Francis de Sales) illustrate the reception, argument, and lived sobriety of the tradition.
\item \textbf{Project synthesis:} The four-term analytical distinction (\emph{dependence}, \emph{duty}, \emph{precept}, \emph{culpability}) summarized in \emph{Conceptual and Legal Distinctions} is an editorial synthesis prepared to keep ontological, moral, juridical, and subjective categories distinct; it carries no independent magisterial or canonical authority.
\end{itemize}

\subsection*{Editions, Versification, and Rights Boundaries}

All verbatim English Scripture quotations are verified against the public-domain Challoner revision of the Douay-Rheims Bible (Project Gutenberg eBook 1581); Latin biblical phrases are verified against the Clementine Vulgate (eBible \texttt{latVUC} verse tables). Psalm numbering follows the Douay-Rheims/Vulgate enumeration, with Hebrew numbering noted alongside at key loci. In particular, Hebrew Psalm~116 is split in the Vulgate and Douay-Rheims into Psalm~114 (Heb.\ 116:1--9) and Psalm~115 (Heb.\ 116:10--19): Challoner retains verse numbers 115:10--19 (used in this essay's citations, e.g., Ps 115:12--13, 16--17), whereas the Clementine Vulgate and the 1899 American Douay-Rheims renumber Psalm~115 as verses 1--10 (making those same verses 115:3--4, 7--8).

Every verbatim quotation from patristic, scholastic, conciliar, or spiritual texts comes from verified public-domain witnesses tracked in the repository source library, together with Emanuel Hoffmann's 1899 public-domain Latin edition of St.\ Augustine's \emph{De civitate Dei} (CSEL 40/1), consulted via the CC~BY-SA~4.0 University of Leipzig / OpenGreekAndLatin TEI encoding. For St.\ Benedict's \emph{Regula Benedicti} (ch.\ 49), the tracked public-domain Latin text is quoted verbatim and accompanied by an original editorial English rendering rather than a published translation. Modern magisterial texts and current canon law (\emph{Sacrosanctum Concilium}, \emph{Lumen gentium}, \emph{Dignitatis humanae}, \emph{Dies Domini}, the \emph{Catechism of the Catholic Church}, the 1983 \emph{Code of Canon Law}, and \emph{Pascite gregem Dei}) are cited solely in accurate doctrinal or legal paraphrase, with zero verbatim reproduction. Historical sixteenth-century catechetical material from the 1566 \emph{Roman Catechism} (such as remarks on civil support of religious feasts) is kept dated within its historical period and never presented as present ecclesiastical law.

\subsection*{Legal Scope and Currentness}

The discussion under \emph{Gift Before Command} includes a canonical illustration of universal Latin Sunday discipline. The following legal bounds control that material:

\begin{itemize}
\item \textbf{Governing body of law and promulgation:} The 1983 \emph{Code of Canon Law} (\emph{Codex Iuris Canonici}), promulgated for the Latin Church by Pope St.\ John Paul~II in the Apostolic Constitution \emph{Sacrae disciplinae leges} (25~January 1983; effective 27~November 1983), together with Book~VI (\emph{De sanctionibus poenalibus in Ecclesia}) as reformed by Pope Francis in the Apostolic Constitution \emph{Pascite gregem Dei} (23~May 2021; effective 8~December 2021).
\item \textbf{Authoritative text, translation status, and citation method:} The official Latin text promulgated in the \emph{Acta Apostolicae Sedis} (and for \emph{Pascite gregem Dei} via \emph{L'Osservatore Romano} and \emph{AAS}) controls. Canons~1, 11, 87 \S 1, 1245, 1246--1248, and revised Book~VI (cc.\ 1311, 1314) are summarized strictly in original plain-English paraphrase without reproducing any published English translation.
\item \textbf{Jurisdiction and persons:} Universal Latin Church only (CIC c.\ 1; c.\ 11). This treatment explicitly excludes the twenty-three Eastern Catholic Churches \emph{sui iuris}, which are governed by the 1990 \emph{Code of Canons of the Eastern Churches} (\emph{Codex Canonum Ecclesiarum Orientalium}, promulgated by \emph{Sacri canones}, 18~October 1990) and their own particular law.
\item \textbf{Material facts assumed:} The summary assumes a baptized or received Catholic of the Latin Church who has completed seven years of age and possesses sufficient use of reason (c.\ 11), with no particular territorial suppression or transfer of holy days, personal dispensation or commutation (cc.\ 87 \S 1, 1245), or grave excusing cause in play.
\item \textbf{Universal and particular law boundary:} Only the universal Latin baseline of cc.\ 87 \S 1, 1245, and 1246--1248 is summarized. Particular calendars of holy days of obligation, suppressions, transfers to Sunday authorized by episcopal conferences with Apostolic See approval, and diocesan or parochial dispensations vary by territory and are outside the scope of this essay.
\item \textbf{Distinction of legal and moral orders:} Canonical obligation (c.\ 1247), fulfillment and impossibility provisions (c.\ 1248), dispensation and commutation authority (cc.\ 87 \S 1, 1245), moral culpability under excusing grave causes, and canonical penalties under revised Book~VI (no \emph{latae sententiae} censure attaches to non-attendance at Sunday Mass) are treated as strictly distinct.
\item \textbf{As-of date, currentness, and check boundary:} Reviewed as of \textbf{2026-11-28} against the cataloged Holy See Latin web delivery states (dated 2026-07-25 for CIC 1983 Books~I--V, VII and revised Book~VI; dated 2026-10-26 for \emph{Sacrae disciplinae leges} and \emph{Pascite gregem Dei}). Printed \emph{Acta Apostolicae Sedis} fascicles, post-delivery universal amendments, and authentic interpretations of the Dicastery (formerly Pontifical Council) for Legislative Texts have not been independently checked beyond those cataloged states.
\item \textbf{Non-advice referral:} These notes serve theological study rather than legal advice, canonical opinion, or pastoral direction. Any individual question concerning dispensations, excusing causes, canonical status, penalties, or sacramental access must be referred to the competent ecclesiastical authority (one's pastor or diocesan bishop) or a qualified canonist furnished with the complete facts.
\end{itemize}

\subsection*{Limitations, Bounded Negative Results, and Review State}

Two bounded source limits are noted openly. First, St.\ Gregory of Nazianzus's \emph{Oration 14} (\emph{On Love for the Poor}) is not included in the tracked public-domain English volume of \emph{Nicene and Post-Nicene Fathers}, Series~II, vol.\ 7, which contains twenty-three selected orations; accordingly, Gregory is cited solely from his tracked \emph{Oration 30.12} on Christ's human will in Gethsemane, and care for the poor is verified directly in St.\ Justin Martyr, St.\ Irenaeus, St.\ John Chrysostom, and St.\ Augustine. Second, Pope St.\ John Paul~II's apostolic letter \emph{Dies Domini} (1998) is not yet registered as a local work in the repository source library and is cited strictly in doctrinal paraphrase for bibliographic orientation alongside tracked patristic and conciliar witnesses.

Remarks on scrupulosity and anxious self-examination under \emph{Failure, Reconciliation, and the Trap of Closing the Ledger} are historical and spiritual observations drawn from St.\ Francis de Sales and Catholic pastoral wisdom; they do not constitute clinical diagnosis, psychological assessment, or medical advice.

Prepared with AI assistance as a private study aid, these pages carry no ecclesiastical \emph{imprimatur}, canonical authority, or official magisterial standing.
